\documentclass[11pt]{article}

\usepackage[utf8]{inputenc}
\usepackage{tabto}
\usepackage[fleqn]{amsmath}
\usepackage{amssymb}
\usepackage{eurosym}
\usepackage{graphicx}
\usepackage{geometry}

\renewcommand{\contentsname}{Inhaltsverzeichnis}

\begin{document}

\begin{titlepage}
	\begin{center}
		\vspace*{1cm}

		\Huge
		\textbf{Tourney}\\
		Analysenotizen

		\vspace{0.5cm}
		\LARGE
		Softwarepraktikum\\
		\Large
		Wintersemester 2014/2015

		\vspace{1.5cm}

		\large
		\textbf{Jonas Auer (2860992)\\
				 Fabian Biester (2859084)\\
				 Jan Tagscherer (2893134)}

		\vfill

		\includegraphics[width=0.4\textwidth]{Logo.png}

		\vspace{1.5cm}

		\Large
		Universität Stuttgart\\
		28.10.2014
	\end{center}
\end{titlepage}

\newpage

\section{Ablauf eines Events}

\begin{itemize}
	\item Beteiligte Personen: Administratoren, Turnierordner (TO), Spieler
	\item Voranmeldung von teilnehmenden Spielern, Registrieren von Bezahlungen
	\item Registrieren von anwesenden Spielern, Vergeben von Startnummern, Zuordnung zu Turnieren auf dem Event
	\item Verteilen der Spieler auf Turniere und Tische mit Paarungen, Festhalten von Ergebnissen mit Laufzetteln
	\item Spielen des Turniers mit den jeweiligen Regeln unter Aufsicht der TOs
	\item Rückmeldung der Ergebnisse an den Administrator und Ausgabe der Abschlusstabelle
\end{itemize}

\section{Ist-Zustand}

\begin{itemize}
	\item Festhalten der vorangemeldeten Spieler und deren Bezahlungsstatus in Excel-Tabellen
	\item Anmeldung der anwesenden Spieler und Vergeben von Startnummern an nur einem Arbeitsplatz möglich, manuelles Eintragen an Hand der Tabelle
	\item Bei Start eines Turniers erneute Anmeldung beim TO nötig
	\item Ausrufen der Paarungen der einzelnen Turniere, TOs bestimmen Paarungen von Hand und halten Turnierergebnisse mit verschiedener Software oder analog fest
	\item Endergebnisse werden an den Administrator übergeben, der diese wiederum händisch übertragen muss
	\item Insgesamt hoher Zeitaufwand und schwere Organisation durch das Ausrufen der Paarungen und die Heterogenität bei der Erfassung von Ergebnissen durch die TOs
\end{itemize}

\newpage
\section{Soll-Zustand}

\begin{itemize}
	\item Modulare Erstellung von Turnierregeln, Cut-Off des Paarungssystems festlegbar (meist Top 16, Top 8 oder Top 4 Spieler)
	\item Anlegen eines Events mit mehreren Turnieren über das Administrationsmodul, Ablegen von Informationen über Turniertyp, Ort und Leitung
	\item Voranmeldung von Spielern, Vermerk des Bezahlungsstatus' und Zuordnung zu Turnieren findet über die Software statt
	\item Persönliche Informationen pro Spieler: Name, Nickname, evtl. Mail-Adresse
	\item Laden der vorangemeldeten Spieler bei Beginn der Veranstaltung und Freischalten bei Anwesenheit oder Anmelden eines neuen Spielers, an mehreren Arbeitsplätzen möglich
	\item Direktes Vergeben von Startnummern, die während des Events bestehen bleiben
	\item Projizieren der stattfindenden Turniere und der jeweiligen Räume über einen Beamer
	\item Übertragen der Daten zu einem Turnier an den jeweiligen TO per Wechseldatenträger
	\item Software legt abhängig von den Spielregeln und eingetragenen Ergebnissen früherer Runden Paarungen fest
	\item Stattfindende Paarungen (Pro Spieler Name, Startnummer und Tischnummer) und Rundentimer können per Beamer projiziert werden
	\item Turnierergebnisse werden nach Ende über einen Wechseldatenträger vom TO wieder zurück zum Administrator übertragen und dort gesammelt
	\item Administrator hat die Möglichkeit, Ergebnistabellen und stattgefundene Paarungen der einzelnen Turniere eines Events auszugeben und auszudrucken
\end{itemize}

\newpage
\section{Weitere Anforderungen an die Software}

\begin{itemize}
	\item Sprache der Software: Primär deutsch, englische Lokalisierung optional
	\item Mengengerüst: Bis zu 128 Spieler pro Turnier realistisch
	\item Mitliefern von Regelmodulen für "Android Netrunner", "Lords of War", "Star Wars: The Card Game", "X-Wing" und "A Game of Thrones"
	\item Von der Software festgelegte Turnierplatzierungen müssen manuell angepasst werden können
	\item Spieler sollten gewarnt werden, wenn sie an zeitlich überlappenden Turnieren teilzunehmen versuchen
	\item Einrichten eines passwortgesicherten Administratorkontos beim ersten Start der Software, über das auch Konten von TOs für Events angelegt und passwortgeschützt werden können
	\item Beim Anlegen eines neuen Turniers sollen entsprechende mögliche Anzahlen von Spielern vorgeschlagen werden
	\item Kommt es zu einem Losverfahren, soll es die Möglichkeiten geben, manuell oder automatisch zu würfeln
	\item Festlegen der angezeigten Informationen bei einer Projektion während den Turnierrunden
	\item Vorschläge bei der Suche von Spielern während der Anmeldung
	\item Vergeben von Freilosen bei ungeraden Spielerzahlen an den schlechtesten Spieler, allerdings wenn möglich nicht mehrfach an gleichen Spieler
	\item Berücksichtigen von Super-Freilosen
	\item Event- und turnierübergreifende Disqualifikation von Spieler möglich, Nachholen von Paarungen im Ermessen des TOs möglich
	\item Warnung bei fehlender Paarung wenn eine neue Runde gestartet wird
	\item Kontrolle der Anwesenheit bei Turnierstart und Verändern der Turnierzuordnungen eines Spielers während des Events möglich
	\item Direktes Wechseln in den Turniermodus bei kleineren Events möglich
	\item Unterstützte Paarungssysteme: Schweizer System, Modifiziertes Schweizer System, K.O.-System, Doppel-K.O.-System, Jeder-gegen-jeden, frei editierbar
\end{itemize}

\end{document}